%% kapa: exam=1 round=30 year=2019 period=1 form=A subject=civil_law subject_ko=민법 seq=19 number=19
%% source: Q-Net 감정평가사 기출문제 (cst00309) artlSeq=5205205 file=extracted/1008/5205205/2019년 제30회 감정평가사 1차시험 문제지/2019년 제30회 감정평가사 1차시험 문제지/2019년도 제30회 감정평가사 1차 1교시 A형.hwp
%% conversion: kordoc 4.24.51 -> tools/gen1.py (auto); review=proofread-sonnet; flags=none
\begin{question}{19}
소멸시효와 제척기간에 관한 설명으로 옳지 않은 것은? (다툼이 있으면 판례에 따름)
\begin{choices}
  \item 소멸시효에 의한 권리소멸은 기산일에 소급하여 효력이 있으나, 제척기간에 의한 권리소멸은 장래에 향하여 효력이 있다.
  \item 소멸시효의 이익은 미리 포기가 가능하나, 제척기간에는 포기가 인정되지 않는다.
  \item 제척기간의 경과는 법원의 직권조사사항이지만, 소멸시효의 완성은 직권조사사항이 아니다.
  \item 소멸시효에는 중단이 인정되고 있으나, 제척기간에는 중단이 인정되지 않는다.
  \item 소멸시효의 정지에 관해서는 민법에 명문의 규정이 있으나, 제척기간의 정지에 관해서는 민법에 명문의 규정이 없다.
\end{choices}
\end{question}
\answer{2}
